%%%%%%%%%%%%%%%%%%%%%%%%%%%%%%%%%%%%%%%
% Cover Letter - Varda Space Industries
% Flight Software Internship, Summer 2027
% Plain XeLaTeX — matches simple resume style (Spectral font)
%%%%%%%%%%%%%%%%%%%%%%%%%%%%%%%%%%%%%%%

\documentclass[11pt]{article}

\usepackage[margin=0.75in, top=0.7in, bottom=0.7in]{geometry}
\usepackage[hidelinks]{hyperref}
\usepackage{parskip}
\usepackage{fontspec}

\setmainfont{Spectral}[
  Path = ../../cv/fonts/,
  UprightFont    = Spectral-Regular.ttf,
  BoldFont       = Spectral-Bold.ttf,
  ItalicFont     = Spectral-Italic.ttf,
  BoldItalicFont = Spectral-BoldItalic.ttf
]

\setlength{\parindent}{0pt}
\setlength{\parskip}{14pt}
\linespread{1.15}
\raggedright

% No hyphenation
\hyphenpenalty=10000
\exhyphenpenalty=10000

\pagestyle{empty}
\urlstyle{same}

\begin{document}

% ── DATE ─────────────────────────────────────────────────────────────────────
\begin{flushright}
  \today
\end{flushright}

% ── BODY ─────────────────────────────────────────────────────────────────────
Dear Varda Hiring Team,

I'm applying for the Flight Software Internship for Summer 2027. I'm a
Computer Science student at RIT in the combined bachelor's and master's
program, with an expected graduation date of December 2027, and have a
strong interest in firmware and embedded software. At GE Vernova, I built a
full I/O subsystem from scratch in C on an STM32F7 for industrial radios.
On the RIT Electric Vehicle Team, I'm the vehicle control unit (VCU)
Software Lead, building real-time ThreadX firmware for our track
motorcycle, including board-level fault handling and health monitoring
across the vehicle's threads and CAN buses. In my free time I'm building an
autonomous fixed-wing flight controller on an ESP32S3, with PID
stabilization, sensor fusion, and GPS waypoint following validated across
an 80-test SITL regression suite; the full airframe recently completed its
first flight in manual mode.

What got me interested in Varda specifically is that the same flight
software has to work across very different phases of one mission, from
quiet on-orbit processing to the heat and deceleration of atmospheric
reentry to landing. What counts as a fault, and how the software should
respond, changes depending on which phase it's in. That's the same kind of
thinking behind the fault handling and health monitoring I built into the
VCU, and I want to bring it to a vehicle that has to work across a mission
like that.

I debug hardware issues regularly with a Saleae logic analyzer, an
oscilloscope, and a J-Link Segger, and I've written drivers directly from
microcontroller datasheets and reference manuals on both the GE Vernova I/O
subsystem and EVT's hardware-agnostic peripheral library. I've worked
across SPI, I2C, UART, ADC, and PWM integrating sensors and actuators into
real-time C and C++
firmware, and the VCU's health-monitoring thread tracks whether every other
thread, the control loop, the CAN receive paths, actually ran correctly
within its expected window. At Lockheed Martin, I also built a code-review
automation tool that became a mandatory part of the team's merge request
process, so I'm comfortable with the testing and process-automation side
of engineering too.

I would welcome the opportunity to talk more about where I could
contribute. Thank you for your consideration.

Thanks,

Travis Brown

\end{document}
